\section*{Bab \ref{ch:b1:stratigraphic-scale} --- Skala stratigrafi}

\begin{solution}{exo:b1:stratigraphic-scale:1}
(a)~Kronostratigrafi, sebuah \omterm{def:b1:stratigraphic-scale:chronostratigraphy}{sistem}. (b)~Litostratigrafi, sebuah \omterm{def:b1:stratigraphic-scale:lithostratigraphy}{formasi}.
(c)~Geokronologi, sebuah epok (batuannya adalah \omterm{def:b1:stratigraphic-scale:chronostratigraphy}{Seri} Kapur Bawah).
(d)~Biostratigrafi, sebuah \omterm{def:b1:stratigraphic-scale:biozone}{biozona}. (e)~Kronostratigrafi, sebuah \omterm{def:b1:stratigraphic-scale:chronostratigraphy}{tahap}.
(f)~Geokronologi, sebuah epok. (g)~Litostratigrafi, sebuah \omterm{def:b1:stratigraphic-scale:lithostratigraphy}{anggota}.
\end{solution}

\begin{solution}{exo:b1:stratigraphic-scale:2}
Paleozoikum $538.8 - 251.902 = \qty{286.9}{\Ma}$, \qty{53.2}{\%} dari
Fanerozoikum; Mesozoikum $251.902 - 66.00 = \qty{185.9}{\Ma}$, \qty{34.5}{\%};
Kenozoikum \qty{66.0}{\Ma}, \qty{12.2}{\%}. Fanerozoikum adalah
$538.8/4567 = \qty{11.8}{\%}$ dari waktu sejak dasar Hadean.
\end{solution}

\begin{solution}{exo:b1:stratigraphic-scale:3}
A, B, C, D diendapkan dalam urutan itu (superposisi), mendatar (horizontalitas asal),
sehingga gradiennya adalah peristiwa kemudian. Dike E memotong A hingga D, sehingga
ia lebih muda daripada D (hubungan memotong); apakah ia datang sebelum atau sesudah
gradien tidak dapat diketahui dari deskripsinya. Permukaan erosi memotong D dan E,
sehingga ia lebih muda daripada keduanya (hubungan memotong). F dan G kemudian
diendapkan di atasnya (superposisi). Sesar menggeser A hingga F tetapi tidak G: ia
bergerak setelah F dan sebelum G (hubungan memotong). Urutan: A, B, C, D; gradien
dan intrusi E (urutan belum ditentukan); erosi; F; sesar; G.
\end{solution}

\begin{solution}{exo:b1:stratigraphic-scale:4}
$r = 30/(85.40 - 84.60) = \qty{37.5}{m}$ per juta tahun. Kemunculan pertama terletak
$12/37.5 = \qty{0.32}{\Ma}$ setelah abu bawah: $85.40 - 0.32 =
\qty{85.08}{\Ma}$.
\end{solution}

\begin{solution}{exo:b1:stratigraphic-scale:5}
(a)~Empat batang vertikal terhadap ketinggian. (b)~B dan C terjadi bersama dari
\qtyrange{26}{50}{m}: sebuah \omterm{def:b1:stratigraphic-scale:biozone}{zona rentang} bersamaan setebal \qty{24}{m}. \omterm{def:b1:stratigraphic-scale:biozone}{Zona interval}
dari kemunculan pertama B hingga kemunculan pertama D membentang dari
\qtyrange{12}{41}{m}, \qty{29}{m}. (c)~Di penampang kedua kemunculan pertama B
(\qty{3}{m}) dan C (\qty{9}{m}) dapat diletakkan, sehingga dasar \omterm{def:b1:stratigraphic-scale:biozone}{zona interval} dari (b)
dan dasar \omterm{def:b1:stratigraphic-scale:biozone}{zona rentang} bersamaan dikenali; puncaknya (kemunculan pertama D) dan puncak
tumpang tindih seperti terlihat di penampang 1 tidak dapat diperiksa, karena D tidak
ada. Kemunculan pertama B dan C ($\qty{12}{m} \leftrightarrow
\qty{3}{m}$, $\qty{26}{m} \leftrightarrow \qty{9}{m}$) adalah garis yang paling aman;
kemunculan terakhir (B pada \qty{50}{m} dan \qty{15}{m}) lebih bergantung pada pengawetan.
\end{solution}

\begin{solution}{exo:b1:stratigraphic-scale:6}
(a)~$60/20 = \qty{3}{\Ma}$: puncak batu pasir tiga juta tahun lebih muda di sumur
sisi darat. (b)~Satu juta tahun lumpur pada \qty{30}{m} per juta tahun: abu terletak
\qty{30}{m} di atas puncak batu pasir. (c)~Ketika abu jatuh, garis pantai telah bergerak
\qty{20}{km} ke pedalaman dari sumur sisi laut, jadi ia masih \qty{40}{km} di sisi laut
dari sumur sisi darat, yang berdiri di darat kering. Abu jatuh di permukaan darat di sana,
di bawah batu pasir yang diendapkan laut dua juta tahun kemudian --- jika ia tidak
tersapu lebih dulu, sebagaimana lapisan lepas di darat mungkin saja tersapu. Lapisan
abu yang sama berada di atas batas \omterm{def:b1:stratigraphic-scale:lithostratigraphy}{formasi} pada satu sumur dan di bawahnya pada sumur
yang lain: batas itu memotong waktu.
\end{solution}

\begin{solution}{exo:b1:stratigraphic-scale:7}
(a)~Lochkovium $1.36/419.62 = \qty{0.32}{\%}$; Induan $0.024/251.902 =
\qty{0.0095}{\%}$, sekitar 34 kali lebih kecil. Dasar Induan terletak di antara lapisan
abu yang ditanggal dengan tepat di penampang GSSP itu sendiri; ketidakpastian yang besar
menunjukkan batas yang umurnya harus dibawa dari level yang bertarikh lebih jauh.
(b)~$251.941 - 251.880 = \qty{0.061}{\Ma} = 61$ ribu tahun, dengan ketidakpastian
$\sqrt{37^2 + 31^2} = 48$ ribu tahun: \num[uncertainty-mode=separate]{60 \pm 48} ribu
tahun yang diterbitkan.
\end{solution}

\begin{solution}{exo:b1:stratigraphic-scale:8}
(a)~$m = 3000 \times \frac43\pi\,(5000)^3 = \qty{1.57e15}{kg}$; debunya,
$60m = \qty{9.42e16}{kg}$. (b)~Permukaan Bumi adalah $4\pi(6.378\times10^6)^2
= \qty{5.11e14}{m^2}$: \qty{184}{kg} per meter persegi. (c)~$184/2500 =
\qty{0.074}{m}$, sekitar \qty{7}{cm}. (d)~Ketebalannya sebanding dengan kubus diameternya:
\qty{1.6}{cm} untuk \qty{6}{km}, \qty{20}{cm} untuk \qty{14}{km}. Orde besarannya, dari
sentimeter hingga desimeter, adalah orde lempung batas, setebal \qty{50}{cm} di El Kef;
perbandingannya hanya dapat kasar, karena debu tidak mungkin jatuh merata dan lempung
juga mengumpulkan sedimen biasa pada zamannya.
\end{solution}

\begin{solution}{exo:b1:stratigraphic-scale:9}
(a)~$\num{60e3} \times \num{2e-6} = \qty{0.12}{m}$, \qty{12}{cm}; untuk 12 dan 108 ribu
tahun, \qty{2.4}{cm} dan \qty{21.6}{cm}. (b)~$\num{32e3} \times
\num{10e-6} = \qty{0.32}{m}$. (c)~Suatu kepunahan menempati sentimeter hingga desimeter
batuan: sampel harus diambil sentimeter demi sentimeter, dan satu sampel setebal
\qty{10}{cm} dalam penampang yang terkondensasi merata-ratakan lima puluh ribu tahun ---
sebagian besar peristiwanya.
\end{solution}

\begin{solution}{exo:b1:stratigraphic-scale:10}
(a)~$u = R/s = 10 \times 2000 = \qty{20000}{m}$ per juta tahun,
\qty{20}{km} per juta tahun; sepanjang \qty{30}{km}, $\Delta t = \qty{1.5}{\Ma}$.
(b)~Dataran kini adalah permukaan $z = s x + a t$. Garis pantai memenuhi
$s x + a t = z_0 + R t$, sehingga $x(t) = (z_0 + (R - a)t)/s$ dan
$\mathrm dx/\mathrm dt = (R - a)/s$, positif (ke pedalaman) selama $a < R$. Ketika
$a > R$ kecepatannya negatif: garis pantai bergerak ke laut, suatu regresi, dan puncak
batu pasir garis pantai kemudian menjadi lebih muda ke arah laut. (c)~Sebuah tebing
sepanjang \qty{200}{m} mencakup, pada \qty{20}{km} per juta tahun, $0.2/20 = \qty{0.01}{\Ma}$
--- sepuluh ribu tahun diakronisme, jauh di bawah apa yang dapat diselesaikan oleh fosil
atau sebagian besar penanggalan; batuan di kedua sisi \omterm{def:b1:geological-maps:geological-map}{kontak} itu sama di sepanjangnya.
\end{solution}

\begin{solution}{exo:b1:stratigraphic-scale:11}
(a)~$f = (25 - 10)/(50 - 10) = 0.375$ dan $t = 100.50 - 0.375 \times 0.80 =
\qty{100.20}{\Ma}$, dengan
\[
  \sigma_t = \sqrt{(0.625 \times 0.05)^2 + (0.375 \times 0.03)^2} = \qty{0.033}{\Ma}.
\] (b)~$g(f) = (1 - f)^2\sigma_1^2 + f^2\sigma_2^2$
memiliki $g'(f) = -2(1 - f)\sigma_1^2 + 2f\sigma_2^2$, nol pada $f^* =
\sigma_1^2/(\sigma_1^2 + \sigma_2^2) = 0.0025/0.0034 = 0.735$, lebih dekat ke lapisan
yang lebih baik penanggalannya; di sana $\sigma_t = \sigma_1\sigma_2/\sqrt{\sigma_1^2 + \sigma_2^2} =
\qty{0.026}{\Ma}$, lebih kecil daripada keduanya. (c)~$g$ adalah parabola yang terbuka
ke atas, sehingga pada $[0, 1]$ ia terbesar di salah satu ujungnya: $g(0) = \sigma_1^2$,
$g(1) = \sigma_2^2$. Maka $\sigma_t \le \max(\sigma_1, \sigma_2)$ di antara kedua lapisan.
\end{solution}

\begin{solution}{exo:b1:stratigraphic-scale:12}
(a)~Batasnya berada di tengah kedalaman: $f = 0.5$, \qty{29.50}{\Ma}. (b)~Dengan
$\varphi_0\lambda = 0.63 \times 1960 = \qty{1234.8}{m}$, ketebalan padat dari
\qtyrange{600}{1000}{m} adalah $400 - 1234.8\,(e^{-600/1960} - e^{-1000/1960}) =
\qty{232.2}{m}$, dan dari \qtyrange{200}{600}{m}, \qty{194.2}{m}; seluruhnya
\qty{426.3}{m}. Fraksi butiran yang diendapkan di antara abu bawah dan batas adalah
$232.2/426.3 = 0.545$, dan umurnya $30.0 - 0.545 =
\qty{29.46}{\Ma}$. (c)~Perkiraan kedua adalah \num{45} ribu tahun lebih muda.
Ia yang lebih baik: separuh yang lebih dalam telah kehilangan lebih banyak airnya, sehingga
ia memuat lebih banyak butiran, dan lebih banyak waktu, per meter daripada separuh yang
lebih dangkal (porositas rerata \num{0.42} berbanding \num{0.51}).
\end{solution}

\subsection*{Masalah: mengorelasikan tiga penampang melintasi sebuah cekungan}

\begin{solution}{pb:b1:stratigraphic-scale:1}
\begin{enumerate}
  \item Tiga kolom pada, misalnya, \qty{1}{cm} untuk \qty{10}{m}, dihubungkan oleh lapisan
        abu dan oleh batas zona; baji Batu Pasir Kell menipis dari \qty{45}{m} di A menjadi
        \qty{6}{m} di C.
  \item Zona~2: A $65.8 - 36.6 = \qty{29.2}{m}$, B \qty{20.8}{m}, C
        \qty{10.0}{m}. Zona~3 yang tercatat: A \qty{14.2}{m}, B \qty{14.0}{m}, C
        \qty{7.2}{m}.
  \item Zona~1 dimulai pada kemunculan pertama \emph{V.~primum}, yang ditemukan dari dasar
        setiap penampang: kemunculan pertamanya terletak di bawah mereka.
  \item Ia hidup hanya di batu pasir: ia mengikuti fasies, dan fasies itu berumur berbeda
        dari tempat ke tempat. Di A ia terdapat di zona~1 dan~2, di B dan C hanya di zona~1.
  \item Di zona~2 di ketiganya (A \qty{49.0}{m}, B \qty{40.0}{m}, C
        \qty{19.0}{m}). Ia konsisten dari penampang ke penampang di sini, dan karena itu dapat
        membagi zona~2, tetapi kemunculan terakhir bergantung pada pengawetan; ia harus
        diperiksa di lebih banyak penampang lebih dulu.
  \item Garis-garisnya menghubungkan \qty{36.6}{m}, \qty{31.2}{m} dan \qty{14.8}{m}, lalu
        \qty{65.8}{m}, \qty{52.0}{m} dan \qty{24.8}{m}. Zona~2 menipis dalam perbandingan
        $29.2 : 20.8 : 10.0$, sekitar $2.9 : 2.1 : 1$: jika zona itu berlangsung sama lama di
        mana-mana, sedimen terakumulasi sekitar tiga kali lebih cepat di A daripada di C.
  \item A: zona~2 ($45.0 > 36.6$); B: zona~1 ($20.0 < 31.2$); C: zona~1
        ($6.0 < 14.8$).
  \item Tidak: ia berada di zona~2 di A, zona~1 di B dan C, sehingga ia lebih muda ke arah
        barat. Serpih laut menyebar ke barat di atas pasir garis pantai: garis pantai bergerak
        ke barat, ke pedalaman, dalam sebuah transgresi.
  \item Sebuah abu jatuh di seluruh cekungan dalam sekejap geologi dan mengendap di atas
        sedimen mana pun; puncak batu pasir mengikuti garis pantai yang bergerak dan bersifat
        \omterm{def:b1:stratigraphic-scale:diachronous}{diakronis}.
  \item Kontinuitas lateral (lapisan itu berlanjut melintasi cekungan), dan superposisi
        (abu~1 di bawah abu~2 di setiap penampang, yang memeriksa korelasinya).
  \item A $(71.0 - 5.0)/1.20 = \qty{55.0}{m}$ per juta tahun; B
        $48.0/1.20 = \qty{40.0}{}$; C $24.0/1.20 = \qty{20.0}{}$.
  \item A paling dekat dengan daratan, tempat sedimen dipasok; lebih sedikit yang sampai ke
        air yang lebih dalam di timur.
  \item $0.1/55 = \num{1.8}$ ribu tahun di A, $\num{2.5}$ di B, $\num{5.0}$
        di C: A menyelesaikan waktu paling halus.
  \item A $152.40 - 40.0/55.0 = \qty{151.67}{\Ma}$; B $152.40 - 12.0/40.0 =
        \qty{152.10}{\Ma}$; C $152.40 - 3.0/20.0 = \qty{152.25}{\Ma}$.
  \item C ke B: $20/(152.25 - 152.10) = \qty{133}{km}$ per juta tahun; B ke A:
        $25/(152.10 - 151.67) = \qty{58.5}{km}$ per juta tahun. Ia melambat.
  \item A $152.40 - 60.8/55.0 = \qty{151.29}{\Ma}$; B $152.40 - 44.0/40.0 =
        \qty{151.30}{\Ma}$; C $152.40 - 21.8/20.0 = \qty{151.31}{\Ma}$.
  \item Ketiganya terentang \num{16} ribu tahun, kurang dari ketidakpastian masing-masing
        (sekitar \num{28} ribu tahun, soal 19): kemunculan pertama itu sinkron di seluruh
        cekungan dalam resolusi penanggalan, dan itulah yang membuatnya menjadi indeks yang baik.
  \item A $152.40 - 31.6/55.0 = \qty{151.826}{\Ma}$; B $152.40 - 23.2/40.0 =
        \qty{151.820}{\Ma}$; C $152.40 - 11.8/20.0 = \qty{151.810}{\Ma}$.
  \item $f = 0.479$, $0.483$, $0.492$; $\sigma_t = \sqrt{(1 - f)^2 + f^2} \times
        0.04 = \qty{0.028}{\Ma}$ di masing-masing, \num{28} ribu tahun.
  \item Rerata \qty{151.819}{\Ma}; sebarannya $151.826 - 151.810 = \num{16}$ ribu
        tahun, lebih kecil daripada ketidakpastian masing-masing umur: ketiganya sepakat.
  \item Interpolasi mengasumsikan satu laju dari abu~1 ke abu~2. Di A interval itu memuat
        batu pasir garis pantai hingga \qty{45}{m} dan serpih di atasnya; pasir dan lumpur tidak
        perlu terakumulasi pada laju yang sama, dan mereka memadat secara berbeda
        (\cref{ch:b1:diagenesis-basins}). Jika pasir terakumulasi lebih cepat, bagian batu pasir
        memuat lebih sedikit waktu daripada yang disarankan ketebalannya, dan kemunculan
        pertama lebih tua daripada yang dihitung --- umur dari A memang yang tertua dari ketiganya.
  \item B: kemunculan pertama terletak di serpih laut, \qty{11}{m} di atas puncak batu pasir,
        tanpa perubahan fasies dekat level itu, dan lajunya memberi \num{2.5} ribu tahun per
        \qty{10}{cm}. A menyelesaikan waktu lebih baik tetapi menempatkan level itu di batu pasir
        garis pantai, \qty{8.4}{m} di bawah perubahan fasies; C terkondensasi, \num{5} ribu tahun
        per \qty{10}{cm}. C yang terburuk untuk resolusi, A untuk fasies.
  \item Dari umur rerata kedua kemunculan pertama, $151.819 - 151.302 =
        \qty{0.52}{\Ma}$: sekitar setengah juta tahun.
  \item Batasnya adalah \qty[uncertainty-mode=separate]{151.82(3)}{\Ma}.
\end{enumerate}
\end{solution}
